\chapter{避孕}\label{ch:g8:contraception}

前面几章描绘了一种力量非凡的机制：
两个\omterm{def:g5:first-look-at-cells:cell}{细胞}相遇，四十周后便有了一个人。
人类是唯一能理解自身机制的\omterm{def:g5:biodiversity:species}{物种}，
这种理解带来了青蛙从不面对的一种选择：
是否以及何时启动这一机制。
\emph{\omterm{prop:g8:contraception:principle}{避孕}}是围绕这一选择的应用科学，
其每种方法都通过打断一个
你现在能够说出的环节来起作用。

\section{原则：说出环节，阻断环节}

\begin{proposition}[避孕与机制相对应]\label{prop:g8:contraception:principle}
\omterm{def:g5:human-reproduction:pregnancy}{怀孕}需要一条不间断的链条
（\cref{prop:g8:sexual-reproduction:core}，
\cref{prop:g8:from-egg-to-birth:firstweek}）：\omterm{def:g8:sexual-reproduction:gametes}{配子}
产生——\omterm{def:g8:reproductive-systems:male}{精子}射入——游向\omterm{def:g8:reproductive-systems:female}{输卵管}
——\omterm{def:g4:animal-reproduction:sexual}{卵细胞}排出并与之相遇——\omterm{def:g5:flower-to-fruit:steps}{受精}——
着床。\emph{避孕}\index{避孕}是
对这条链条的任何有意打断。药房货架上的每种方法
都是对某个已命名环节的打击，而
按所针对的环节进行分类，正是整个理论的核心：
\begin{itemize}
  \item \emph{屏障}类方法阻止\omterm{def:g8:sexual-reproduction:gametes}{配子}相遇；
  \item \emph{\omterm{def:g8:puberty:hormones}{激素}}类方法根本阻止\omterm{def:g4:animal-reproduction:sexual}{卵细胞}
        被排出；
  \item 其他方法则使相遇或着床
        失败。
\end{itemize}
\end{proposition}
\admitted

\section{主要方法，按其环节分类}

\begin{definition}[屏障类方法：避孕套]\label{def:g8:contraception:condom}
\emph{避孕套}\index{避孕套}——一种薄鞘，戴于
\omterm{def:g8:reproductive-systems:male}{阴茎}上（女性版本衬于阴道内）——阻断
链条的射精环节：精液永远不会进入
伴侣体内。随处可得，仅在需要时使用
——且在货架上独具第二项作用：它是
唯一能阻挡
经性接触传播的\emph{感染}，其内容
在\cref{ch:g9:microbes-and-infection}中讲述。这一双重用途
正是\omterm{prop:g8:contraception:principle}{避孕}套即使在用别种方法时仍保留的
缘由。
\end{definition}

\begin{definition}[激素类方法：避孕药]\label{def:g8:contraception:pill}
\emph{避孕药}\index{避孕药}
利用\cref{def:g8:puberty:hormones}中的同一机制：每日
服用，其\omterm{def:g8:puberty:hormones}{激素}使\omterm{def:g8:reproductive-systems:female}{卵巢}的周期
处于停滞状态
（\cref{prop:g8:reproductive-systems:cycle}中的事件 4，
即无\omterm{def:g5:human-reproduction:pregnancy}{怀孕}发生）——\emph{不\omterm{prop:g8:reproductive-systems:cycle}{排卵}，无\omterm{def:g4:animal-reproduction:sexual}{卵细胞}，
无可\omterm{def:g5:flower-to-fruit:steps}{受精}之物}。由医生开具，医生
据人选方法；规律服用时可靠性很高
——却对感染完全无效，
因此常须与\omterm{def:g8:contraception:condom}{避孕套}长期配合。
\end{definition}

\begin{example}[货架上的方法，按环节分类]\label{ex:g8:contraception:shelf}
其他常见方法，按其所针对的环节归档：
\omterm{def:g8:puberty:hormones}{激素}埋植剂与贴片——原理同\omterm{def:g8:contraception:pill}{避孕药}，
但持续时间更长（数月或数年，无须每日服用）；
医生置于\omterm{prop:g5:human-reproduction:plan}{子宫}内的小型装置——作用于
\omterm{def:g5:flower-to-fruit:steps}{受精}与着床环节，使其失效
长达数年；以及\emph{紧急\omterm{def:g8:contraception:pill}{避孕药}}，是数日内
服用的事后\omterm{def:g8:puberty:hormones}{激素}剂量，主要
通过推迟尚未发生的\omterm{prop:g8:reproductive-systems:cycle}{排卵}起效——是一种
最后手段，可在药店获取，不能
替代\emph{事先}使用的方法。以上这些方法都
无法防范感染。
\end{example}

\begin{remark}[无效的方法]\label{rem:g8:contraception:notwork}
操场上流传的方法，对照这条链条来看：
“数安全期”败在
\cref{rem:g8:reproductive-systems:living}所说的变异性
——\omterm{prop:g8:reproductive-systems:cycle}{排卵}日期无法提前得知，而
\omterm{def:g8:reproductive-systems:male}{精子}可在其间存活数日；“及时退出”
败在\omterm{def:g8:reproductive-systems:male}{精子}早在射精之前就已游出的事实。
两者失败之处，恰与按环节命名的理论所预测一致：
二者都无法可靠地阻断任何环节。传言并不是方法。
\end{remark}

\section{选择、责任与求助指南}

\begin{proposition}[选择是真实的，也是共同的]\label{prop:g8:contraception:choice}
机制本身所决定的事实：
\begin{itemize}
  \item \omterm{def:g5:human-reproduction:pregnancy}{怀孕}只需一次行为，不取决于关系的长短：
        链条不计次数；
  \item 这是一条\emph{共同}的链条：双方的\omterm{def:g5:first-look-at-cells:cell}{细胞}，
        双方的责任——\omterm{def:g8:contraception:condom}{避孕套}与
        \omterm{def:g8:contraception:pill}{避孕药}都不只是某一方的事，而是
        一场沟通；
  \item 除禁欲外，没有哪种方法是完美的，不过
        好的方法在严格按设计使用时已十分接近
        完美——“正确使用”恰恰是
        现实中多数失败的所在。
\end{itemize}
\end{proposition}
\admitted

\begin{example}[求助指南]\label{ex:g8:contraception:help}
\Cref{ex:g8:puberty:questions}中的求助渠道，再作补充：
学校护士与家庭医生可为\omterm{prop:g8:contraception:principle}{避孕}问题
提供保密咨询，包括未成年人；
计划生育服务中心在每个地区都有，正是
为此而设，免费且不加评判；药剂师负责
处理紧急\omterm{def:g8:contraception:pill}{避孕药}的时限。体系的设计，使得
无人须凭传言去摸索这条链条——求助
指南本身也是方法的一部分。
\end{example}

\begin{method}[评估你听说的任何方法]\label{met:g8:contraception:evaluate}
对任何声称有效的方法，依次追问：
\begin{enumerate}
  \item 它阻断链条中的\emph{哪一环节}——
        能否在图上指出它？
  \item 按真实人群的使用方式，\emph{可靠性}如何？
  \item 它是否也能预防\emph{感染}，
        抑或须由\omterm{def:g8:contraception:condom}{避孕套}来配合？
  \item 它需要什么——每日记忆、
        一张处方、一次试配——这些是否
        适合使用者？
\end{enumerate}
无法指明的环节，便不成其为方法
（\cref{rem:g8:contraception:notwork}
在第 1 问即不通过）。
\end{method}

\section{练习}

\begin{exercise}[$\star$]\label{exo:g8:contraception:1}
写出\omterm{def:g5:human-reproduction:pregnancy}{怀孕}所需的链条，并
据此给出\omterm{prop:g8:contraception:principle}{避孕}的定义。
\end{exercise}

\begin{exercise}[$\star$]\label{exo:g8:contraception:2}
\omterm{def:g8:contraception:condom}{避孕套}阻断了哪一环节？\omterm{def:g8:contraception:pill}{避孕药}呢？
\end{exercise}

\begin{exercise}[$\star$]\label{exo:g8:contraception:3}
单凭\omterm{def:g8:contraception:condom}{避孕套}，还能提供哪项另外的服务？
\end{exercise}

\begin{exercise}[$\star$]\label{exo:g8:contraception:4}
紧急\omterm{def:g8:contraception:pill}{避孕药}主要通过何种方式起作用，
它不能做什么？
\end{exercise}

\begin{exercise}[$\star$]\label{exo:g8:contraception:5}
写出求助图上的三个去处，
并说明各自提供什么服务。
\end{exercise}

\begin{exercise}[$\star$]\label{exo:g8:contraception:6}
为什么“数安全期”不可靠？需给出两条理由，
依据\cref{rem:g8:contraception:notwork}。
\end{exercise}

\begin{exercise}[$\star\star$]\label{exo:g8:contraception:7}
将\cref{ex:g8:contraception:shelf}中的货架
按所攻击的环节分类，并加入\omterm{def:g8:contraception:condom}{避孕套}和\omterm{def:g8:contraception:pill}{避孕药}。
\end{exercise}

\begin{exercise}[$\star\star$]\label{exo:g8:contraception:8}
医生为何常对年轻伴侣说：“\omterm{def:g8:contraception:pill}{避孕药}\emph{和}
\omterm{def:g8:contraception:condom}{避孕套}”？两种防护，各用一句话说明。
\end{exercise}

\begin{exercise}[$\star\star$]\label{exo:g8:contraception:9}
借助
\cref{prop:g8:reproductive-systems:cycle}，说明\omterm{def:g8:contraception:pill}{避孕药}将周期保持在何种状态，
并解释这为何能防止
\omterm{def:g5:human-reproduction:pregnancy}{怀孕}。
\end{exercise}

\begin{exercise}[$\star\star$]\label{exo:g8:contraception:10}
“双方的\omterm{def:g5:first-look-at-cells:cell}{细胞}，双方的事。”
请为这句话辩护，反驳把\omterm{prop:g8:contraception:principle}{避孕}交给伴侣
一方的习惯。
\end{exercise}

\begin{exercise}[$\star\star$]\label{exo:g8:contraception:11}
依据\cref{met:g8:contraception:evaluate}，分析：(a) \omterm{def:g8:puberty:hormones}{激素}
植入物；(b) “及时退出”。
\end{exercise}

\begin{exercise}[$\star\star\star$]\label{exo:g8:contraception:12}
“\omterm{prop:g8:contraception:principle}{避孕}是应用生理学：货架上的每种方法，
都是本书的一章，各有所指。”请以三种方法
及其对应的三章加以论证。
\end{exercise}

\section{问题：计划生育咨询}

\begin{problem}[{周末问题——计划生育中心
一个下午的问题}]\label{pb:g8:contraception:1}
一位计划生育咨询师（虚构）的一个下午：
四位来访者，四种情境，墙上
挂着一张链条图。

\medskip\noindent\textbf{第一部分——墙上的
链条图。}
\begin{enumerate}
  \item 画出墙上的链条图：从\omterm{def:g8:sexual-reproduction:gametes}{配子}到
        着床的完整链条，并标出
        \cref{prop:g8:contraception:principle}中的三个攻击点。
  \item 来访者 1 问：“一天一片药，怎么能
        阻止任何事？”请从机制回答：
        药物用的是哪种信号语言，又向谁发出？
  \item 来访者 1 追问：“那要是我忘了
        吃呢？”请在图上找出危险所在——
        \omterm{def:g8:puberty:hormones}{激素}信号中断后，什么会\omterm{def:g7:organs-at-work:recovery}{恢复}？
  \item 咨询师为何无论选择哪种方法，
        都仍推荐\omterm{def:g8:contraception:condom}{避孕套}？
        （双重职责。）
\end{enumerate}

\medskip\noindent\textbf{第二部分——各种情境。}
\begin{enumerate}[resume]
  \item 来访者 2 今年十七岁，焦虑不安，
        相信一款“安全期”应用让其他方法变得不再必要。
        请温和地纠正这一想法：
        哪两条生物学事实会
        使日历法失效？
  \item 来访者 3 需要一种“无需每日
        记忆”的方法。请提供货架上的两种候选，
        并说明它们分别阻断的环节。
  \item 来访者 4 在\omterm{def:g8:contraception:condom}{避孕套}破裂后一天内
        前来就诊。请说出可选方案、其时限、
        其主要机制——以及咨询师事后将讨论的
        长期方法。
  \item 四位来访者离开时，都带着同一句
        关于感染的结束语。请写出这句话。
\end{enumerate}

\medskip\noindent\textbf{第三部分——总结。}
\begin{enumerate}[resume]
  \item 实习生问，中心为何教授链条，
        而不只是发放方法？
        请借助
        \cref{met:g8:contraception:evaluate}中的
        第一个问题来作答。
  \item 将这一下午涉及的方法，按各自对
        使用者的要求分类（依据方法中的第 4 问）。
  \item 实习生不解，为何未成年人的就诊
        在设计上即为保密？请借助
        \cref{ex:g8:contraception:help}中的思路作答。
  \item 用一句话总结这一下午：中心
        除药盒之外，实际提供的是什么。
\end{enumerate}
\end{problem}
